% Experiment 8: backend/evals/paper/exp8_scale.py   (run alone; five repetitions)
% Language channels: backend/evals/run_eval.py --repeats 5
\emph{Cost.} Table~\ref{tab:scale} times the evidence pipeline on synthetic campuses
of 1,200 to 10,000 students with about 20 payments each, against one Neo4j 5.26
instance. Each graph job was run five times on the loaded graph and we report the
median with the range; each trust lookup is one call of the production
\code{social\_scores} for five candidates, timed over 50 calls after 10 discarded
warm-up calls. The experiment ran alone on the machine described in
Section~\ref{sec:eval-method}. At 10,000 students the three graph jobs together took
11.9~s, against a scheduler that runs them every ten minutes, and the trust lookup on
the dispatch path took 3.0~ms at the median and 4.2~ms at the 95th percentile. The
ring search now returns every ring, 117 of them on the largest synthetic campus,
whose random payments close many cycles; this is what the earlier cap of 50 hid. The
production Tarjan routine processed a million nodes and three million edges in
9.4~s, and fitting the acceptance model took 2.5~s on 10,000 logged rounds and 25~s
on 100,000, about a semester's rounds for a campus of 5,000 students.

\begin{table*}[!t]
\caption{Cost of the Evidence Pipeline. Median (Range) of Five Repetitions}
\label{tab:scale}
\centering
\footnotesize
\setlength{\tabcolsep}{6pt}
\begin{tabular}{@{}rrrcccc@{}}
\toprule
& & & \multicolumn{3}{c}{Graph jobs (s)} & Trust lookup \\
\cmidrule(lr){4-6}
Students & Friendships & Payments & Closure & Circulation & Rings & p50 / p95 (ms) \\
\midrule
1,200 & 4,050 & 23,944 & 0.11 (0.06--0.14) & 0.81 (0.63--0.93) & 0.51 (0.50--0.54) & 5.7 / 14.8 \\
2,500 & 8,408 & 49,895 & 0.14 (0.12--1.16) & 1.65 (1.57--2.90) & 1.25 (1.05--1.41) & 10.9 / 16.5 \\
5,000 & 16,862 & 99,787 & 0.32 (0.26--0.40) & 3.19 (3.04--3.91) & 2.41 (2.20--2.63) & 4.0 / 6.3 \\
10,000 & 33,755 & 199,525 & 0.54 (0.49--0.88) & 6.53 (6.13--7.27) & 4.91 (4.59--4.95) & 3.0 / 4.2 \\
\midrule
\multicolumn{7}{@{}l}{Tarjan routine alone, 3 out-edges per node: $10^4$ nodes 0.03~s;
$10^5$ nodes 0.66~s; $10^6$ nodes 9.4~s} \\
\multicolumn{7}{@{}l}{Acceptance-model fit: $10^3$ rounds 0.24~s; $10^4$ rounds 2.5~s;
$10^5$ rounds 25~s} \\
\bottomrule
\multicolumn{7}{@{}p{0.8\textwidth}@{}}{\scriptsize Hardware: Intel Core i7-12650H, 16 threads,
7.6~GiB available to Docker; Python 3.12; Neo4j 5.26 Community. Trust lookups: 50
timed calls after 10 warm-up calls, five candidates each.}
\end{tabular}
\end{table*}

\emph{Language channels.} The advisory channels play a small part in this design,
and we tested them only for that part. On eleven hand-written cases with five runs
each, the errand-text channel scored genuine groups 41.7 points higher than farmed
ones on a 0--100 scale and cleared every honest control, but it also cleared one of
three farmed groups whose errands had been written to look varied. The
rating-comment channel showed no separation ($-1.2$ points) and failed to clear an
honest runner whose friends rate without comments, so it is switched off. Neither
channel followed instructions injected into the text. The cases were written by the
authors, and the harness refuses to report an accuracy figure until twenty
blind-labelled pilot cases exist, so these numbers show only that the errand-text
channel works end to end, not that it is accurate.

\emph{Tests.} The full suite of 273 automated tests, including at least one
regression test for each of the six faults in Section~\ref{sec:impl}, passed in
111~s.
